% GENERATED FILE. Edit canonical lessons, then run npm run generate:books.
% canonical-source-hash: fnv1a64:06e5a6f1596f2764
% canonical-lessons: PA-C260-sava, PA-C260-sadhe, PA-C260-paune, PA-C260-dedh, PA-C260-dhai

\chapter{Quarter Past and Half Past}
\label{ch:pa-a2-quarter-past-and-half-past}
\hlchaptermodality{260}

\begin{chapteropening}
\textbf{By the end of this chapter:} \emph{I can say a time with a quarter or a half.}

Complete the last lesson of chapter 260: I can say a time with a quarter or a half.
\end{chapteropening}

\section[\texorpdfstring{\pa{ਸਵਾ} (savā)}{ਸਵਾ (savā)}]{\pa{ਸਵਾ} (savā) — plus a quarter; a quarter past}

\label{lesson:PA-C260-sava}



\begin{quote}
Before the new one: say the Punjabi for a voice, then the Punjabi for to dial a number.
\end{quote}

\subsection*{You'll want to know: \pa{ਸਵਾ}}
\textbf{\pa{ਸਵਾ}} — \emph{savā} — \textquotedblleft{}plus a quarter; a quarter past\textquotedblright{}.

\begin{quote}
\pa{ਬੱਸ} \pa{ਸਵਾ} \pa{ਅੱਠ} \pa{ਵਜੇ} \pa{ਆਉਂਦੀ} \pa{ਹੈ।}
\end{quote}

\begin{quote}
\emph{bass savā aṭṭh vaje āundī hai.}
\end{quote}

\begin{quote}
The bus comes at a quarter past eight.
\end{quote}

\begin{cousinweb}
\textbf{\pa{ਸਵਾ}} (\emph{savā}) is inherited from Sanskrit \emph{sapāda}, \textquotedblleft{}increased by a fourth\textquotedblright{}; CDIAL lists Punjabi \emph{savā}, \textquotedblleft{}plus a quarter\textquotedblright{}, under it (CDIAL 13134).
\end{cousinweb}

\begin{grammarlens}[title={What you've built}]
One more word for reading notices and messages.
\end{grammarlens}

\subsection*{Your turn}
\begin{itemize}
\raggedright
  \item \emph{Say it:} \emph{savā}
  \item \emph{Say it:} \emph{savā}, once more
  \item \emph{Say it:} \emph{bass savā aṭṭh vaje āundī hai.} — The bus comes at a quarter past eight.
  \item \emph{Recall:} say the Punjabi for a voice, then the Punjabi for to dial a number, then say \emph{savā} again
\end{itemize}

\begin{tcolorbox}[breakable,colback=teal!4,colframe=teal!35!black,title={Before you move on}]
What does \pa{ਸਵਾ} mean? (\textquotedblleft{}Plus a quarter; a quarter past\textquotedblright{}.) Say it once more.
\end{tcolorbox}
\section[\texorpdfstring{\pa{ਸਾਢੇ} (sāḍhe)}{ਸਾਢੇ (sāḍhe)}]{\pa{ਸਾਢੇ} (sāḍhe) — plus a half; half past}

\label{lesson:PA-C260-sadhe}



\begin{quote}
Before the new one: say the Punjabi for to dial a number, then the Punjabi for plus a quarter.
\end{quote}

\subsection*{You'll want to know: \pa{ਸਾਢੇ}}
\textbf{\pa{ਸਾਢੇ}} — \emph{sāḍhe} — \textquotedblleft{}plus a half; half past\textquotedblright{}.

\begin{quote}
\pa{ਦਫ਼ਤਰ} \pa{ਦਾ} \pa{ਕੰਮ} \pa{ਸਾਢੇ} \pa{ਨੌਂ} \pa{ਵਜੇ} \pa{ਸ਼ੁਰੂ} \pa{ਹੁੰਦਾ} \pa{ਹੈ।}
\end{quote}

\begin{quote}
\emph{daftar dā kamm sāḍhe nauṁ vaje shurū hundā hai.}
\end{quote}

\begin{quote}
Work at the office starts at half past nine.
\end{quote}

\begin{cousinweb}
\textbf{\pa{ਸਾਢੇ}} (\emph{sāḍhe}) is inherited from Sanskrit \emph{sārdha}, \textquotedblleft{}plus a half\textquotedblright{}; CDIAL lists Punjabi \emph{sāḍhe}, \textquotedblleft{}plus a half\textquotedblright{}, under it (CDIAL 13369).
\end{cousinweb}

\begin{grammarlens}[title={What you've built}]
One more word for reading notices and messages.
\end{grammarlens}

\subsection*{Your turn}
\begin{itemize}
\raggedright
  \item \emph{Say it:} \emph{sāḍhe}
  \item \emph{Say it:} \emph{sāḍhe}, once more
  \item \emph{Say it:} \emph{daftar dā kamm sāḍhe nauṁ vaje shurū hundā hai.} — Work at the office starts at half past nine.
  \item \emph{Recall:} say the Punjabi for to dial a number, then the Punjabi for plus a quarter, then say \emph{sāḍhe} again
  \item \emph{Write it:} copy \textbf{\pa{ਸਾਢੇ}} once, checking the whole word
\end{itemize}

\begin{tcolorbox}[breakable,colback=teal!4,colframe=teal!35!black,title={Before you move on}]
What does \pa{ਸਾਢੇ} mean? (\textquotedblleft{}Plus a half; half past\textquotedblright{}.) Say it once more.
\end{tcolorbox}
\section[\texorpdfstring{\pa{ਪੌਣੇ} (pauṇe)}{ਪੌਣੇ (pauṇe)}]{\pa{ਪੌਣੇ} (pauṇe) — less a quarter; a quarter to}

\label{lesson:PA-C260-paune}



\begin{quote}
Before the new one: say the Punjabi for plus a quarter, then the Punjabi for plus a half.
\end{quote}

\subsection*{You'll want to know: \pa{ਪੌਣੇ}}
\textbf{\pa{ਪੌਣੇ}} — \emph{pauṇe} — \textquotedblleft{}less a quarter; a quarter to\textquotedblright{}.

\begin{quote}
\pa{ਰੇਲ} \pa{ਪੌਣੇ} \pa{ਦਸ} \pa{ਵਜੇ} \pa{ਜਾਂਦੀ} \pa{ਹੈ।}
\end{quote}

\begin{quote}
\emph{rel pauṇe das vaje jāndī hai.}
\end{quote}

\begin{quote}
The train goes at a quarter to ten.
\end{quote}

\begin{cousinweb}
\textbf{\pa{ਪੌਣੇ}} (\emph{pauṇe}) is inherited from a reconstructed \emph{*padūna}, \textquotedblleft{}less by a quarter\textquotedblright{}; CDIAL lists Punjabi \emph{pauṇe}, \textquotedblleft{}three quarters\textquotedblright{}, under it (CDIAL 7765).
\end{cousinweb}

\begin{grammarlens}[title={What you've built}]
One more word for reading notices and messages.
\end{grammarlens}

\subsection*{Your turn}
\begin{itemize}
\raggedright
  \item \emph{Say it:} \emph{pauṇe}
  \item \emph{Say it:} \emph{pauṇe}, once more
  \item \emph{Say it:} \emph{rel pauṇe das vaje jāndī hai.} — The train goes at a quarter to ten.
  \item \emph{Recall:} say the Punjabi for plus a quarter, then the Punjabi for plus a half, then say \emph{pauṇe} again
  \item \emph{Write it:} copy \textbf{\pa{ਪੌਣੇ}} once, checking the whole word
\end{itemize}

\begin{tcolorbox}[breakable,colback=teal!4,colframe=teal!35!black,title={Before you move on}]
What does \pa{ਪੌਣੇ} mean? (\textquotedblleft{}Less a quarter; a quarter to\textquotedblright{}.) Say it once more.
\end{tcolorbox}
\section[\texorpdfstring{\pa{ਡੇਢ} (ḍeḍh)}{ਡੇਢ (ḍeḍh)}]{\pa{ਡੇਢ} (ḍeḍh) — one and a half}

\label{lesson:PA-C260-dedh}



\begin{quote}
Before the new one: say the Punjabi for plus a half, then the Punjabi for less a quarter.
\end{quote}

\subsection*{You'll want to know: \pa{ਡੇਢ}}
\textbf{\pa{ਡੇਢ}} — \emph{ḍeḍh} — \textquotedblleft{}one and a half\textquotedblright{}.

\begin{quote}
\pa{ਸਟੇਸ਼ਨ} \pa{ਘਰ} \pa{ਤੋਂ} \pa{ਡੇਢ} \pa{ਘੰਟਾ} \pa{ਦੂਰ} \pa{ਹੈ।}
\end{quote}

\begin{quote}
\emph{saṭeshan ghar ton ḍeḍh ghaṇṭā dūr hai.}
\end{quote}

\begin{quote}
The station is an hour and a half from home.
\end{quote}

\begin{cousinweb}
\textbf{\pa{ਡੇਢ}} (\emph{ḍeḍh}) is inherited from Sanskrit \emph{dvyardha}, \textquotedblleft{}one and a half\textquotedblright{}; CDIAL lists Punjabi \emph{ḍeḍh}, \textquotedblleft{}one and a half\textquotedblright{}, under it (CDIAL 6698).
\end{cousinweb}

\begin{grammarlens}[title={What you've built}]
One more word for reading notices and messages.
\end{grammarlens}

\subsection*{Your turn}
\begin{itemize}
\raggedright
  \item \emph{Say it:} \emph{ḍeḍh}
  \item \emph{Say it:} \emph{ḍeḍh}, once more
  \item \emph{Say it:} \emph{saṭeshan ghar ton ḍeḍh ghaṇṭā dūr hai.} — The station is an hour and a half from home.
  \item \emph{Recall:} say the Punjabi for plus a half, then the Punjabi for less a quarter, then say \emph{ḍeḍh} again
  \item \emph{Write it:} copy \textbf{\pa{ਡੇਢ}} once, checking the whole word
\end{itemize}

\begin{tcolorbox}[breakable,colback=teal!4,colframe=teal!35!black,title={Before you move on}]
What does \pa{ਡੇਢ} mean? (\textquotedblleft{}One and a half\textquotedblright{}.) Say it once more.
\end{tcolorbox}
\section[\texorpdfstring{\pa{ਢਾਈ} (ḍhāī)}{ਢਾਈ (ḍhāī)}]{\pa{ਢਾਈ} (ḍhāī) — two and a half}

\label{lesson:PA-C260-dhai}



\begin{quote}
Before the new one: say the Punjabi for less a quarter, then the Punjabi for one and a half.
\end{quote}

\subsection*{You'll want to know: \pa{ਢਾਈ}}
\textbf{\pa{ਢਾਈ}} — \emph{ḍhāī} — \textquotedblleft{}two and a half\textquotedblright{}.

\begin{quote}
\pa{ਮੈਂ} \pa{ਢਾਈ} \pa{ਵਜੇ} \pa{ਘਰ} \pa{ਆਵਾਂਗਾ।}
\end{quote}

\begin{quote}
\emph{maiṁ ḍhāī vaje ghar āvāṅgā.}
\end{quote}

\begin{quote}
I will come home at half past two.
\end{quote}

\begin{cousinweb}
\textbf{\pa{ਢਾਈ}} (\emph{ḍhāī}) is inherited from Sanskrit \emph{ardhatṛtīya}, \textquotedblleft{}two and a half\textquotedblright{}; CDIAL lists Punjabi \emph{ḍhāī}, \textquotedblleft{}two and a half\textquotedblright{}, under it (CDIAL 651).
\end{cousinweb}

\begin{grammarlens}[title={What you've built}]
One more word for reading notices and messages.
\end{grammarlens}

\subsection*{Your turn}
\begin{itemize}
\raggedright
  \item \emph{Say it:} \emph{ḍhāī}
  \item \emph{Say it:} \emph{ḍhāī}, once more
  \item \emph{Say it:} \emph{maiṁ ḍhāī vaje ghar āvāṅgā.} — I will come home at half past two.
  \item \emph{Recall:} say the Punjabi for less a quarter, then the Punjabi for one and a half, then say \emph{ḍhāī} again
  \item \emph{Write it:} copy \textbf{\pa{ਢਾਈ}} once, checking the whole word
\end{itemize}

\begin{tcolorbox}[breakable,colback=teal!4,colframe=teal!35!black,title={Before you move on}]
What does \pa{ਢਾਈ} mean? (\textquotedblleft{}Two and a half\textquotedblright{}.) Say it once more.
\end{tcolorbox}
